\documentclass[10pt]{article}
% Standard packages
\usepackage{fullpage}
\usepackage[normalem]{ulem} % For strikethrough font

\usepackage{graphicx}

\usepackage{pgfkeys}

% Required packages for tikz-uml
\usepackage{tikz}
\usetikzlibrary{positioning}
\usepackage{ifthen}
\usepackage{xstring}
\usepackage{calc}
\usepackage{pgfplots}
\usepackage{tikz-uml}
\newcommand{\dunder}{\rule{0.1in}{0.1pt}}
\newcommand{\Newpage}{\end{preview}\begin{preview}}
\newcommand{\setvalue}[1]{\pgfkeys{/variables/#1}}
\newcommand{\getvalue}[1]{\pgfkeysvalueof{/variables/#1}}
\newcommand{\declare}[1]{%
 \pgfkeys{
  /variables/#1.is family,
  /variables/#1.unknown/.style = {\pgfkeyscurrentpath/\pgfkeyscurrentname/.initial = ##1}
 }%
}

\title{Robotic Path Planning Software Design Documentation}
% Name
\author{Bruun, E. \and Flores, J. \and Kumar, V.}

%Just to make the uml class code work here without modification
\declare{} 
\newcommand{\relstr}{} 

\begin{document}
\maketitle

\section{Introduction (EB \& JF)}
This software is designed to solve the classic problem of robotic path planning. Given a start and end point, a path must be found that does not collide with a set of obstacles. The space for which the path can be found is not discretized, so it does not have existing nodes for the robot to travel between. Due to this, a graph of vertices and edges must be created. In addition to generating the graph, determining which is the shortest path must also be solved. 

Currently, the path planning problem is one of the most researched topics in autonomous robotics. There are many path planning algorithms that address this problem. Many algorithms employ a grid based search, but we aim to solve the path planning problem with an interval based search to cover the entirety of the free space instead of a grid.  

This software focuses on implementing Rapidly-exploring random trees (RRT) algorithms. RRT is a common option that both creates a graph and finds a path, but the path will not necessarily be optimal. RRT*, popularized by Dr. Karaman and Dr. Frazzoli, is an optimized modified algorithm that aims to achieve a shortest path, whether by distance or other metrics. Although these are the algorithms implemented, other solvers could be added and utilized in the future.

This software will be developed such that it can be implemented in both 2D and 3D. The ultimate goal is to utilize this software for the design of structures such as bridges that span a certain distance while avoiding specified obstacles in the way. The intended audience is structural designers working on preliminary designs. The actual structural analysis of the resulting structure is out of the scope of this project, but could easily be added as a feature in the future. 

\newpage
\section{Program Classes}
Show the division and main components of the program

\subsection{Input (VK)}
Reading in info and defining domain

\begin{tikzpicture}
\begin{umlpackage}{Input Classes}
	%% Input and Parameter Classes %%

%Parameters (ABSTRACT)
\setvalue{Xcoord=16.}
\setvalue{Ycoord=2.}
\setvalue{aligninfo=north}
\setvalue{widthinfo=6cm}
\renewcommand{\relstr}{}
%This class provides a standard interface to an input file which provides at run-time for program parameters such as step_size
\umlclass[x=\getvalue{Xcoord}, y=\getvalue{Ycoord},
				type = abstract,
				\relstr,
				anchor=\getvalue{aligninfo},
				width = \getvalue{widthinfo}]
				{Params}
% Data
{
+ \textit{file$\_$name: string}
}
% Methods
{
+\textit{\dunder{init}\dunder(filename)}\\
+ \textit{parse$\_$parameters(string file$\_$name)} \\
+ void declare$\_$entry(string entry$\_$name, \\ string default$\_$value, \\ Reg$\_$Ex(?) pattern)\\
+ \textit{set(string entry$\_$name, string new$\_$value)}\\
+ \textit{print$\_$paramters()}
}
	

%Domain Parameters
\setvalue{aligninfo=north}
\setvalue{widthinfo=5cm}
\renewcommand{\relstr}{below right= 0cm and 9cm of Params.north}
%This class provides a standard interface to an input file which provides at run-time for program parameters such as step_size
\umlclass[x=\getvalue{Xcoord}, y=\getvalue{Ycoord},
				\relstr,
				anchor=\getvalue{aligninfo},
				width = \getvalue{widthinfo}]
				{ParamsDomain}
% Data
{
+ file$\_$name: string \\
+ dims : int \\
+ domain : list[name str, info] \\ %name str = type of geometry (i.e. point, circle, rectangle, freeform)
+ obstacles : list[name str, info] %name str = type of geometry (i.e. point, circle, rectangle, freeform)
}
% Methods
{
}
	
\umlinherit[anchors= 180 and 0, geometry=-|-]{ParamsDomain}{Params}

%Origin/Goal Parameters
\setvalue{aligninfo=north}
\setvalue{widthinfo=5cm}
\renewcommand{\relstr}{below right=3cm and 9cm of Params.north}
%This class provides a standard interface to an input file which provides at run-time for program parameters such as step_size
\umlclass[x=\getvalue{Xcoord}, y=\getvalue{Ycoord},
				\relstr,
				anchor = \getvalue{aligninfo},
				width = \getvalue{widthinfo}]
				{ParamsOriginGoal}
% Data
{
+ file$\_$name : string \\
+ origin : list[name str, info]\\ %name str = type of geometry (i.e. point, circle, rectangle, freeform)
+ goal : list[name str, info]\\ %name str = type of geometry (i.e. point, circle, rectangle, freeform)
}
% Methods
{
}
	
\umlinherit[anchors= 180 and 0, geometry=-|-]{ParamsOriginGoal}{Params}

%RRT Parameters
\setvalue{aligninfo=north}
\setvalue{widthinfo=5cm}
\renewcommand{\relstr}{below right=6cm and 9cm of Params.north}
%This class provides a standard interface to an input file which provides at run-time for program parameters such as step_size
\umlclass[x=\getvalue{Xcoord}, y=\getvalue{Ycoord},
				\relstr,
				anchor=\getvalue{aligninfo},
				width = \getvalue{widthinfo}]
				{ParamsRRT}
% Data
{
+ file$\_$name : string \\
+ method : string \\
+ trials : int\\
+ step$\_$size : float\\
+ additional$\_$params :% Placeholder for any necessary params
}
% Methods
{
}
	
\umlinherit[anchors= 180 and 0, geometry=-|-]{ParamsRRT}{Params}

\end{umlpackage}

\end{tikzpicture}

\subsection{Geometry (VK \& JF)}
Shapes used for domain, obstacles

\subsection{Domain (VK \& EB)}
Assembling the domain in which algorithm acts

\subsection{Algorithm (EB)}
RRT implemetation

\subsection{Output (JF)}
Output of the program


\newpage
\section{UML Diagram and Dependencies}
\includegraphics[trim={1.5cm 2cm 3cm 5cm},clip, width=\columnwidth]{../uml_class_diagram/project_class_diagram.pdf}

\subsection{Discussion}

\newpage
\section{Algorithms}

\subsection{RRT}
Add description and pseudocode 
\subsection{RRT*}
Add description and pseudocode 

\subsection{Testing}
State the approach to verify the functional correctness of the modules.

\section{Results}
\subsection{RRT}
\subsection{RRT*}

\section{Profiling}
Speed up attempt, what was slow etc. 

\section{Conclusion}
Comparison and conclusions, how it was working together

\section{References}


\end{document}
